\chapter*{运行环境与操作说明}
\addcontentsline{toc}{chapter}{运行环境与操作说明}
软件启动后，通过设备栏连接 Nexys A7-100T 开发板。实时监测、事件记录、听歌识曲和体感游戏共用该设备连接，主要操作如下。

{\small\noindent\begin{tabularx}{\linewidth}{p{27mm}Y}
\toprule 步骤 & 操作与可观察结果\\\midrule
1．下载启动 & 点击封面的 EXE 链接。连接开发板的 PROG/UART USB 接口并上电，在设备栏选择实际设备。\\
2．准备开发板 & 首次使用或掉电后，确认板型并点击“初始化此板”，将内置配套固件加载到 SRAM。已有配套程序时直接连接。\\
3．实时感知 & 查看声音、频谱、XYZ、温度、电压与开关按键；轻倾开发板，观察 XYZ 变化。点击“抓取1.5秒”，查看完整波形并播放录音。\\
4．智能识曲 & 手机靠近板载麦克风播放歌曲；SW15 拨上或点击进入识曲，等待就绪后连续敲桌三下。观察录音、传输、匹配与结果；可中止任务、查看历史并收藏。\\
5．体感交互 & SW15 拨下并退出识曲，进入“体感游戏”，完成握持与方向校准。轻倾移动、BTNC 中键刹车；无板时选择键盘练习。\\
\bottomrule
\end{tabularx}}

\begin{tcolorbox}[title=版本与运行条件]
正式软件：1.6.0；配套固件：\code{ec0351f8}；目标器件：\code{xc7a100tcsg324-1}。适用 Windows 10/11 64 位。软件内置运行环境、固件与下载后端；首次缺少 USB 驱动时按软件引导准备。初始化不写 Flash，开发板掉电后需重新初始化。在线识曲需要网络，本地指纹和键盘练习可在本机运行。
\end{tcolorbox}

\begin{cnabstract}
基于 Nexys A7-100T 开发板，设计并实现了多模态感知与交互系统。FPGA采集板载麦克风、三轴加速度、温度、片内电压及开关按键信号，使用统一时间戳组织数据；主机软件完成实时显示、声振事件记录、听歌识曲和体感游戏控制。

硬件采用 SystemVerilog 实现PDM/CIC音频转换、SPI与I²C接口、XADC监测、特征提取、双组事件缓存及UART传输。事件窗口包含触发前0.5秒和触发后1秒；音乐录音通过31抽头FIR低通滤波和二倍抽取，在160000点缓冲内保存20秒、8 kHz音频。主机采用单串口采集线程、GUI消息分发和独立识曲进程，在线识曲使用ShazamIO，本地备用采用音频指纹检索，支持三敲启动、任务中止、历史查询和收藏。体感游戏将校准后的XYZ映射为移动输入，以中键实现刹车，并与识曲及自动事件记录进行互斥管理。

固件在100 MHz约束下完成布线，建立时间最差裕量为0.102 ns，保持时间最差裕量为0.026 ns，BRAM占用119.5/135，DSP占用9/240。系统通过100项集成检查、完整事件导出、配置恢复和EXE资源校验；三首板载麦克风录音的在线识曲回放均匹配成功，体感游戏完成八关实板通关测试。
\end{cnabstract}
\cnkeywords{FPGA；多模态感知；软硬件协同；音频指纹；体感交互}
